\documentclass{amsart}
% For dealing with references we use the comment environment
\usepackage{verbatim}
\newenvironment{reference}{\comment}{\endcomment}
%\newenvironment{reference}{}{}
\newenvironment{slogan}{\comment}{\endcomment}
\newenvironment{history}{\comment}{\endcomment}

% For commutative diagrams we use Xy-pic
\usepackage[all]{xy}

% We use 2cell for 2-commutative diagrams.
\xyoption{2cell}
\UseAllTwocells

% We use multicol for the list of chapters between chapters
\usepackage{multicol}

% This is generally recommended for better output
\usepackage{lmodern}
\usepackage[T1]{fontenc}

% For cross-file-references

% Package for hypertext links:
\usepackage{hyperref}

% For any local file, say "hello.tex" you want to link to please

% Theorem environments.
%
\theoremstyle{plain}
\newtheorem{theorem}[subsection]{Theorem}
\newtheorem{proposition}[subsection]{Proposition}
\newtheorem{lemma}[subsection]{Lemma}

\theoremstyle{definition}
\newtheorem{definition}[subsection]{Definition}
\newtheorem{example}[subsection]{Example}
\newtheorem{exercise}[subsection]{Exercise}
\newtheorem{situation}[subsection]{Situation}

\theoremstyle{remark}
\newtheorem{remark}[subsection]{Remark}
\newtheorem{remarks}[subsection]{Remarks}

\numberwithin{equation}{subsection}

% Macros
%
\def\lim{\mathop{\mathrm{lim}}\nolimits}
\def\colim{\mathop{\mathrm{colim}}\nolimits}
\def\Spec{\mathop{\mathrm{Spec}}}
\def\Hom{\mathop{\mathrm{Hom}}\nolimits}
\def\Ext{\mathop{\mathrm{Ext}}\nolimits}
\def\SheafHom{\mathop{\mathcal{H}\!\mathit{om}}\nolimits}
\def\SheafExt{\mathop{\mathcal{E}\!\mathit{xt}}\nolimits}
\def\Sch{\mathit{Sch}}
\def\Mor{\mathop{\mathrm{Mor}}\nolimits}
\def\Ob{\mathop{\mathrm{Ob}}\nolimits}
\def\Sh{\mathop{\mathit{Sh}}\nolimits}
\def\NL{\mathop{N\!L}\nolimits}
\def\CH{\mathop{\mathrm{CH}}\nolimits}
\def\proetale{{pro\text{-}\acute{e}tale}}
\def\etale{{\acute{e}tale}}
\def\QCoh{\mathit{QCoh}}
\def\Ker{\mathop{\mathrm{Ker}}}
\def\Im{\mathop{\mathrm{Im}}}
\def\Coker{\mathop{\mathrm{Coker}}}
\def\Coim{\mathop{\mathrm{Coim}}}

% Boxtimes
%
\DeclareMathSymbol{\boxtimes}{\mathbin}{AMSa}{"02}

%
% Macros for moduli stacks/spaces
%
\def\QCohstack{\mathcal{QC}\!\mathit{oh}}
\def\Cohstack{\mathcal{C}\!\mathit{oh}}
\def\Spacesstack{\mathcal{S}\!\mathit{paces}}
\def\Quotfunctor{\mathrm{Quot}}
\def\Hilbfunctor{\mathrm{Hilb}}
\def\Curvesstack{\mathcal{C}\!\mathit{urves}}
\def\Polarizedstack{\mathcal{P}\!\mathit{olarized}}
\def\Complexesstack{\mathcal{C}\!\mathit{omplexes}}
% \Pic is the operator that assigns to X its picard group, usage \Pic(X)
% \Picardstack_{X/B} denotes the Picard stack of X over B
% \Picardfunctor_{X/B} denotes the Picard functor of X over B
\def\Pic{\mathop{\mathrm{Pic}}\nolimits}
\def\Picardstack{\mathcal{P}\!\mathit{ic}}
\def\Picardfunctor{\mathrm{Pic}}
\def\Deformationcategory{\mathcal{D}\!\mathit{ef}}

\setlength{\emergencystretch}{3em}
\begin{document}
\title[Stacks Project, Tag 0F8E]{2112 --- A bounded below quasi-coherent complex representing a perfect object on a quasi-compact quasi-separated scheme with the resolution property admits a quasi-isomorphism from a bounded complex of finite locally free modules}
\maketitle
\noindent Copyright (C) 2005--2025 Johan de Jong. Stacks Project authors. Source: \href{https://stacks.math.columbia.edu/tag/0F8E}{Tag 0F8E}.

\noindent Editorial difficulty note (not author text). Difficulty H1: The given complex may consist of arbitrary quasi-coherent modules and the scheme need not be Noetherian or have an ample line bundle. The author constructs a global finite-vector-bundle complex by selecting finite-type cycle submodules, resolving them through the resolution property, and reducing the Tor-amplitude through a cone induction. This is a global strictification construction beyond ordinary local projective-resolution calculations..

\noindent Editorial provenance note (not author text). History: excerpt from revision \texttt{a0444\allowbreak 6e57e\allowbreak c1fbc\allowbreak 252a8\allowbreak 71afc\allowbreak ec775\allowbreak 2fb28\allowbreak 07b14}, perfect.tex, lines 9572--9670. Reference presentation was adapted; original-author context and core support blocks were retained or added. The mathematical wording is unchanged. Licensed under GNU FDL 1.2 or later; no invariant sections or cover texts. See ../licenses/COPYING. Context blocks reproduce author definitions and supporting results; they are not additional counted problems.

\begin{definition}
\label{definition-resolution-property}
Let $X$ be a scheme. We say $X$ has the {\it resolution property}
if every quasi-coherent $\mathcal{O}_X$-module of finite type
is the quotient of a finite locally free $\mathcal{O}_X$-module.
\end{definition}

\begin{definition}
\label{cohomology-definition-perfect}
Let $(X, \mathcal{O}_X)$ be a ringed space.
Let $\mathcal{E}^\bullet$ be a complex of $\mathcal{O}_X$-modules.
We say $\mathcal{E}^\bullet$ is {\it perfect} if there exists
an open covering $X = \bigcup U_i$ such that for each $i$
there exists a morphism of complexes
$\mathcal{E}_i^\bullet \to \mathcal{E}^\bullet|_{U_i}$
which is a quasi-isomorphism with $\mathcal{E}_i^\bullet$
a strictly perfect complex of $\mathcal{O}_{U_i}$-modules.
An object $E$ of $D(\mathcal{O}_X)$ is {\it perfect}
if it can be represented by a perfect complex of $\mathcal{O}_X$-modules.
\end{definition}

\begin{definition}
\label{cohomology-definition-strictly-perfect}
Let $(X, \mathcal{O}_X)$ be a ringed space.
Let $\mathcal{E}^\bullet$ be a complex of $\mathcal{O}_X$-modules.
We say $\mathcal{E}^\bullet$ is {\it strictly perfect}
if $\mathcal{E}^i$ is zero for all but finitely many $i$ and
$\mathcal{E}^i$ is a direct summand of a finite free
$\mathcal{O}_X$-module for all $i$.
\end{definition}

\begin{definition}
\label{cohomology-definition-tor-amplitude}
Let $(X, \mathcal{O}_X)$ be a ringed space.
Let $E$ be an object of $D(\mathcal{O}_X)$.
Let $a, b \in \mathbf{Z}$ with $a \leq b$.
\begin{enumerate}
\item We say $E$ has {\it tor-amplitude in $[a, b]$}
if $H^i(E \otimes_{\mathcal{O}_X}^\mathbf{L} \mathcal{F}) = 0$
for all $\mathcal{O}_X$-modules $\mathcal{F}$ and all $i \not \in [a, b]$.
\item We say $E$ has {\it finite tor dimension}
if it has tor-amplitude in $[a, b]$ for some $a, b$.
\item We say $E$ {\it locally has finite tor dimension}
if there exists an open covering $X = \bigcup U_i$ such that
$E|_{U_i}$ has finite tor dimension for all $i$.
\end{enumerate}
An $\mathcal{O}_X$-module $\mathcal{F}$ has {\it tor dimension $\leq d$}
if $\mathcal{F}[0]$ viewed as an object of $D(\mathcal{O}_X)$ has
tor-amplitude in $[-d, 0]$.
\end{definition}

\begin{definition}
\label{cohomology-definition-pseudo-coherent}
Let $(X, \mathcal{O}_X)$ be a ringed space. Let $\mathcal{E}^\bullet$
be a complex of $\mathcal{O}_X$-modules. Let $m \in \mathbf{Z}$.
\begin{enumerate}
\item We say $\mathcal{E}^\bullet$ is {\it $m$-pseudo-coherent}
if there exists an open covering $X = \bigcup U_i$ and for each $i$
a morphism of complexes
$\alpha_i : \mathcal{E}_i^\bullet \to \mathcal{E}^\bullet|_{U_i}$
where $\mathcal{E}_i^\bullet$ is strictly perfect on $U_i$ and
$H^j(\alpha_i)$ is an isomorphism for $j > m$ and $H^m(\alpha_i)$
is surjective.
\item We say $\mathcal{E}^\bullet$ is {\it pseudo-coherent}
if it is $m$-pseudo-coherent for all $m$.
\item We say an object $E$ of $D(\mathcal{O}_X)$ is
{\it $m$-pseudo-coherent} (resp.\ {\it pseudo-coherent})
if and only if it can be represented by a $m$-pseudo-coherent
(resp.\ pseudo-coherent) complex of $\mathcal{O}_X$-modules.
\end{enumerate}
\end{definition}

\begin{definition}
\label{derived-definition-cone}
Let $\mathcal{A}$ be an additive category.
Let $f : K^\bullet \to L^\bullet$ be a morphism of
complexes of $\mathcal{A}$. The {\it cone} of $f$
is the complex $C(f)^\bullet$ given by
$C(f)^n = L^n \oplus K^{n + 1}$ and
differential
$$
d_{C(f)}^n =
\left(
\begin{matrix}
d^n_L & f^{n + 1} \\
0 & -d_K^{n + 1}
\end{matrix}
\right)
$$
It comes equipped with canonical morphisms of complexes
$i : L^\bullet \to C(f)^\bullet$ and $p : C(f)^\bullet \to K^\bullet[1]$
induced by the obvious maps $L^n \to C(f)^n \to K^{n + 1}$.
\end{definition}

\begin{definition}
\label{homology-definition-quasi-isomorphism-cochain}
Let $\mathcal{A}$ be an abelian category.
\begin{enumerate}
\item A morphism of cochain complexes $f : A^\bullet \to B^\bullet$
of $\mathcal{A}$ is called a {\it quasi-isomorphism} if the induced
maps $H^i(f) : H^i(A^\bullet) \to H^i(B^\bullet)$
is an isomorphism for all $i \in \mathbf{Z}$.
\item A cochain complex $A^\bullet$ is called
{\it acyclic} if all of its cohomology objects
$H^i(A^\bullet)$ are zero.
\end{enumerate}
\end{definition}

\begin{definition}
\label{derived-definition-split-ses}
Let $\mathcal{A}$ be an additive category.
A {\it termwise split exact sequence of complexes of $\mathcal{A}$}
is a complex of complexes
$$
0 \to
A^\bullet \xrightarrow{\alpha}
B^\bullet \xrightarrow{\beta}
C^\bullet \to 0
$$
together with given direct sum decompositions
$B^n = A^n \oplus C^n$
compatible with $\alpha^n$ and $\beta^n$.
We often write $s^n : C^n \to B^n$ and $\pi^n : B^n \to A^n$
for the maps induced by the direct sum decompositions.
According to
Homology, Lemma \href{https://stacks.math.columbia.edu/tag/011J}{lemma ses termwise split cochain (Tag 011J)}
we get an associated morphism of complexes
$$
\delta : C^\bullet \longrightarrow A^\bullet[1]
$$
which in degree $n$ is the map $\pi^{n + 1} \circ d_B^n \circ s^n$.
In other words
$(A^\bullet, B^\bullet, C^\bullet, \alpha, \beta, \delta)$
forms a triangle
$$
A^\bullet \to B^\bullet \to C^\bullet \to A^\bullet[1]
$$
This will be the {\it triangle associated to the termwise
split sequence of complexes}.
\end{definition}

\begin{lemma}
\label{properties-lemma-quasi-coherent-colimit-finite-type}
Let $X$ be a quasi-compact and quasi-separated scheme.
Any quasi-coherent sheaf of $\mathcal{O}_X$-modules
is the directed colimit of its quasi-coherent
$\mathcal{O}_X$-submodules which are of finite type.
\end{lemma}

\begin{proof}
The colimit is directed because if $\mathcal{G}_1$, $\mathcal{G}_2$
are quasi-coherent subsheaves of finite type, then the image of
$\mathcal{G}_1 \oplus \mathcal{G}_2 \to \mathcal{F}$ is
a quasi-coherent submodule of finite type.
Let $U \subset X$ be any affine open, and let
$s \in \Gamma(U, \mathcal{F})$ be any section.
Let $\mathcal{G} \subset \mathcal{F}|_U$ be the
subsheaf generated by $s$. Then clearly $\mathcal{G}$
is quasi-coherent and has finite type as an $\mathcal{O}_U$-module.
By Lemma \ref{properties-lemma-extend} we see that $\mathcal{G}$ is the restriction
of a quasi-coherent subsheaf $\mathcal{G}' \subset \mathcal{F}$
which has finite type. Since $X$ has a basis for the topology consisting
of affine opens we conclude that every local section of
$\mathcal{F}$ is locally contained in a quasi-coherent submodule
of finite type. Thus we win.
\end{proof}


\begin{lemma}
\label{properties-lemma-extend}
Let $X$ be a quasi-compact and quasi-separated scheme.
Let $U \subset X$ be a quasi-compact open.
Let $\mathcal{F}$ be a quasi-coherent $\mathcal{O}_X$-module.
Let $\mathcal{G} \subset \mathcal{F}|_U$ be a quasi-coherent
$\mathcal{O}_U$-submodule which is of finite type. Then
there exists a quasi-coherent submodule $\mathcal{G}' \subset \mathcal{F}$
which is of finite type such that $\mathcal{G}'|_U = \mathcal{G}$.
\end{lemma}

\begin{proof}
Let $n$ be the minimal number of affine opens $U_i \subset X$,
$i = 1, \ldots , n$ such that $X = U \cup \bigcup U_i$.
(Here we use that $X$ is quasi-compact.) Suppose
we can prove the lemma for the case $n = 1$. Then we can successively
extend $\mathcal{G}$
to a $\mathcal{G}_1$ over $U \cup U_1$
to a $\mathcal{G}_2$ over $U \cup U_1 \cup U_2$
to a $\mathcal{G}_3$ over $U \cup U_1 \cup U_2 \cup U_3$,
and so on.
Thus we reduce to the case $n = 1$.

\medskip\noindent
Thus we may assume that $X = U \cup V$ with $V$ affine.
Since $X$ is quasi-separated and $U$, $V$ are quasi-compact open,
we see that $U \cap V$ is a quasi-compact open. It suffices to prove the
lemma for the system $(V, U \cap V, \mathcal{F}|_V, \mathcal{G}|_{U \cap V})$
since we can glue the resulting sheaf $\mathcal{G}'$ over $V$
to the given sheaf $\mathcal{G}$ over $U$ along the common value
over $U \cap V$.
Thus we reduce to the case where $X$ is affine.

\medskip\noindent
Assume $X = \Spec(R)$. Write $\mathcal{F} = \widetilde M$
for some $R$-module $M$. By Lemma \ref{properties-lemma-extend-trivial} above we may
find a quasi-coherent subsheaf $\mathcal{H} \subset \mathcal{F}$
which restricts to $\mathcal{G}$ over $U$.
Write $\mathcal{H} = \widetilde N$ for some $R$-module $N$.
For every $u \in U$ there exists an $f \in R$ such that
$u \in D(f) \subset U$ and such that $N_f$ is finitely generated,
see Lemma \href{https://stacks.math.columbia.edu/tag/01PB}{lemma finite type module (Tag 01PB)}.
Since $U$ is quasi-compact we can cover it by finitely
many $D(f_i)$ such that $N_{f_i}$ is generated by
finitely many elements, say $x_{i, 1}/f_i^N, \ldots, x_{i, r_i}/f_i^N$.
Let $N' \subset N$ be the submodule generated by the elements
$x_{i, j}$. Then the subsheaf
$\mathcal{G}' = \widetilde{N'} \subset \mathcal{H} \subset \mathcal{F}$
works.
\end{proof}


\begin{lemma}
\label{properties-lemma-extend-trivial}
Let $j : U \to X$ be a quasi-compact open immersion of schemes.
\begin{enumerate}
\item Any quasi-coherent sheaf on $U$ extends to a quasi-coherent
sheaf on $X$.
\item Let $\mathcal{F}$ be a quasi-coherent sheaf on $X$.
Let $\mathcal{G} \subset \mathcal{F}|_U$ be a quasi-coherent
subsheaf. There exists a quasi-coherent subsheaf $\mathcal{H}$ of
$\mathcal{F}$ such that $\mathcal{H}|_U = \mathcal{G}$
as subsheaves of $\mathcal{F}|_U$.
\item Let $\mathcal{F}$ be a quasi-coherent sheaf on $X$.
Let $\mathcal{G}$ be a quasi-coherent sheaf on $U$.
Let $\varphi : \mathcal{G} \to \mathcal{F}|_U$ be a morphism
of $\mathcal{O}_U$-modules. There exists a quasi-coherent sheaf $\mathcal{H}$
of $\mathcal{O}_X$-modules and a map $\psi : \mathcal{H} \to \mathcal{F}$
such that $\mathcal{H}|_U = \mathcal{G}$ and that
$\psi|_U = \varphi$.
\end{enumerate}
\end{lemma}

\begin{proof}
An immersion is separated
(see Schemes, Lemma \href{https://stacks.math.columbia.edu/tag/01L7}{lemma immersions monomorphisms (Tag 01L7)})
and $j$ is quasi-compact by assumption.
Hence for any quasi-coherent sheaf $\mathcal{G}$ on $U$ the sheaf
$j_*\mathcal{G}$ is an extension to $X$. See
Schemes, Lemma \href{https://stacks.math.columbia.edu/tag/01LC}{lemma push forward quasi coherent (Tag 01LC)} and
Sheaves, Section \href{https://stacks.math.columbia.edu/tag/009Z}{section open immersions (Tag 009Z)}.

\medskip\noindent
Assume $\mathcal{F}$, $\mathcal{G}$ are as in (2).
Then $j_*\mathcal{G}$ is a quasi-coherent sheaf on $X$ (see above).
It is a subsheaf of $j_*j^*\mathcal{F}$.
Hence the kernel
$$
\mathcal{H} =
\Ker(\mathcal{F} \oplus j_* \mathcal{G}
\longrightarrow j_*j^*\mathcal{F})
$$
is quasi-coherent as well, see
Schemes, Section \href{https://stacks.math.columbia.edu/tag/01LA}{section quasi coherent (Tag 01LA)}.
It is formal to check that $\mathcal{H} \subset \mathcal{F}$ and that
$\mathcal{H}|_U = \mathcal{G}$ (using the material in
Sheaves, Section \href{https://stacks.math.columbia.edu/tag/009Z}{section open immersions (Tag 009Z)} again).

\medskip\noindent
Part (3) is proved in the same manner as (2). Just take
$\mathcal{H} = \Ker(\mathcal{F} \oplus j_* \mathcal{G}
\to j_*j^*\mathcal{F})$ with its obvious map to $\mathcal{F}$
and its obvious identification with $\mathcal{G}$ over $U$.
\end{proof}


\begin{lemma}
\label{cohomology-lemma-finite-cohomology}
Let $(X, \mathcal{O}_X)$ be a ringed space.
Let $K$ be an object of $D(\mathcal{O}_X)$.
Let $m \in \mathbf{Z}$.
\begin{enumerate}
\item If $K$ is $m$-pseudo-coherent and $H^i(K) = 0$
for $i > m$, then $H^m(K)$ is a finite type $\mathcal{O}_X$-module.
\item If $K$ is $m$-pseudo-coherent and $H^i(K) = 0$
for $i > m + 1$, then $H^{m + 1}(K)$ is a finitely presented
$\mathcal{O}_X$-module.
\end{enumerate}
\end{lemma}

\begin{proof}
Proof of (1). We may work locally on $X$. Hence we may assume there exists
a strictly perfect complex $\mathcal{E}^\bullet$ and a map
$\alpha : \mathcal{E}^\bullet \to K$ which induces
an isomorphism on cohomology in degrees $> m$ and a surjection in degree $m$.
It suffices to prove the result for $\mathcal{E}^\bullet$.
Let $n$ be the largest integer such that $\mathcal{E}^n \not = 0$.
If $n = m$, then $H^m(\mathcal{E}^\bullet)$ is a quotient of
$\mathcal{E}^n$ and the result is clear.
If $n > m$, then $\mathcal{E}^{n - 1} \to \mathcal{E}^n$ is surjective as
$H^n(E^\bullet) = 0$. By Lemma \ref{cohomology-lemma-local-lift-map}
we can locally find a section of this surjection and write
$\mathcal{E}^{n - 1} = \mathcal{E}' \oplus \mathcal{E}^n$.
Hence it suffices to prove the result for the complex
$(\mathcal{E}')^\bullet$ which is the same as $\mathcal{E}^\bullet$
except has $\mathcal{E}'$ in degree $n - 1$ and $0$ in degree $n$.
We win by induction on $n$.

\medskip\noindent
Proof of (2). We may work locally on $X$. Hence we may assume there exists
a strictly perfect complex $\mathcal{E}^\bullet$ and a map
$\alpha : \mathcal{E}^\bullet \to K$ which induces
an isomorphism on cohomology in degrees $> m$ and a surjection in degree $m$.
As in the proof of (1) we can reduce to the case that $\mathcal{E}^i = 0$
for $i > m + 1$. Then we see that
$H^{m + 1}(K) \cong H^{m + 1}(\mathcal{E}^\bullet) =
\Coker(\mathcal{E}^m \to \mathcal{E}^{m + 1})$
which is of finite presentation.
\end{proof}


\begin{lemma}
\label{cohomology-lemma-local-lift-map}
Let $(X, \mathcal{O}_X)$ be a ringed space.
Given a solid diagram of $\mathcal{O}_X$-modules
$$
\xymatrix{
\mathcal{E} \ar@{..>}[dr] \ar[r] & \mathcal{F} \\
& \mathcal{G} \ar[u]_p
}
$$
with $\mathcal{E}$ a direct summand of a finite free
$\mathcal{O}_X$-module and $p$ surjective, then a dotted arrow
making the diagram commute exists locally on $X$.
\end{lemma}

\begin{proof}
We may assume $\mathcal{E} = \mathcal{O}_X^{\oplus n}$ for some $n$.
In this case finding the dotted arrow is equivalent to lifting the
images of the basis elements in $\Gamma(X, \mathcal{F})$. This is
locally possible by the characterization of surjective maps of
sheaves (Sheaves, Section \href{https://stacks.math.columbia.edu/tag/007S}{section exactness points (Tag 007S)}).
\end{proof}


\begin{lemma}
\label{cohomology-lemma-perfect}
Let $(X, \mathcal{O}_X)$ be a ringed space.
Let $E$ be an object of $D(\mathcal{O}_X)$.
The following are equivalent
\begin{enumerate}
\item $E$ is perfect, and
\item $E$ is pseudo-coherent and locally has finite tor dimension.
\end{enumerate}
\end{lemma}

\begin{proof}
Assume (1). By definition this means there exists an open covering
$X = \bigcup U_i$ such that $E|_{U_i}$ is represented by a
strictly perfect complex. Thus $E$ is pseudo-coherent (i.e.,
$m$-pseudo-coherent for all $m$) by
Lemma \href{https://stacks.math.columbia.edu/tag/08CC}{lemma pseudo coherent independent representative (Tag 08CC)}.
Moreover, a direct summand of a finite free module is flat, hence
$E|_{U_i}$ has finite Tor dimension by
Lemma \ref{cohomology-lemma-tor-amplitude}. Thus (2) holds.

\medskip\noindent
Assume (2). After replacing $X$ by the members of an open covering
we may assume there exist integers $a \leq b$ such that $E$
has tor amplitude in $[a, b]$. Since $E$ is $m$-pseudo-coherent
for all $m$ we conclude using Lemma \ref{cohomology-lemma-perfect-precise}.
\end{proof}


\begin{lemma}
\label{cohomology-lemma-perfect-precise}
Let $(X, \mathcal{O}_X)$ be a ringed space.
Let $E$ be an object of $D(\mathcal{O}_X)$.
Let $a \leq b$ be integers. If $E$ has tor amplitude in $[a, b]$
and is $(a - 1)$-pseudo-coherent, then $E$ is perfect.
\end{lemma}

\begin{proof}
After replacing $X$ by the members of an open covering we may assume there
exists a strictly perfect complex $\mathcal{E}^\bullet$ and a map
$\alpha : \mathcal{E}^\bullet \to E$ such that $H^i(\alpha)$ is an isomorphism
for $i \geq a$. We may and do replace $\mathcal{E}^\bullet$ by
$\sigma_{\geq a - 1}\mathcal{E}^\bullet$. Choose a distinguished triangle
$$
\mathcal{E}^\bullet \to E \to C \to \mathcal{E}^\bullet[1]
$$
From the vanishing of cohomology sheaves of $E$ and $\mathcal{E}^\bullet$
and the assumption on $\alpha$ we obtain $C \cong \mathcal{K}[2 - a]$ with
$\mathcal{K} = \Ker(\mathcal{E}^{a - 1} \to \mathcal{E}^a)$.
Let $\mathcal{F}$ be an $\mathcal{O}_X$-module.
Applying $- \otimes_{\mathcal{O}_X}^\mathbf{L} \mathcal{F}$
the assumption that $E$ has tor amplitude in $[a, b]$
implies $\mathcal{K} \otimes_{\mathcal{O}_X} \mathcal{F} \to
\mathcal{E}^{a - 1} \otimes_{\mathcal{O}_X} \mathcal{F}$ has image
$\Ker(\mathcal{E}^{a - 1} \otimes_{\mathcal{O}_X} \mathcal{F}
\to \mathcal{E}^a \otimes_{\mathcal{O}_X} \mathcal{F})$.
It follows that $\text{Tor}_1^{\mathcal{O}_X}(\mathcal{E}', \mathcal{F}) = 0$
where $\mathcal{E}' = \Coker(\mathcal{E}^{a - 1} \to \mathcal{E}^a)$.
Hence $\mathcal{E}'$ is flat (Lemma \href{https://stacks.math.columbia.edu/tag/08BQ}{lemma flat tor zero (Tag 08BQ)}).
Thus $\mathcal{E}'$ is locally a direct summand of a finite free module by
Modules, Lemma \href{https://stacks.math.columbia.edu/tag/08BL}{lemma flat locally finite presentation (Tag 08BL)}.
Thus locally the complex
$$
\mathcal{E}' \to \mathcal{E}^{a + 1} \to \ldots \to \mathcal{E}^b
$$
is quasi-isomorphic to $E$ and $E$ is perfect.
\end{proof}


\begin{lemma}
\label{cohomology-lemma-tor-amplitude}
Let $(X, \mathcal{O}_X)$ be a ringed space.
Let $E$ be an object of $D(\mathcal{O}_X)$.
Let $a, b \in \mathbf{Z}$ with $a \leq b$. The following are equivalent
\begin{enumerate}
\item $E$ has tor-amplitude in $[a, b]$.
\item $E$ is represented by a complex
$\mathcal{E}^\bullet$ of flat $\mathcal{O}_X$-modules with
$\mathcal{E}^i = 0$ for $i \not \in [a, b]$.
\end{enumerate}
\end{lemma}

\begin{proof}
If (2) holds, then we may compute
$E \otimes_{\mathcal{O}_X}^\mathbf{L} \mathcal{F} =
\mathcal{E}^\bullet \otimes_{\mathcal{O}_X} \mathcal{F}$
and it is clear that (1) holds.

\medskip\noindent
Assume that (1) holds. We may represent $E$ by a bounded above complex
of flat $\mathcal{O}_X$-modules $\mathcal{K}^\bullet$, see
Section \href{https://stacks.math.columbia.edu/tag/06Y7}{section flat (Tag 06Y7)}.
Let $n$ be the largest integer such that $\mathcal{K}^n \not = 0$.
If $n > b$, then $\mathcal{K}^{n - 1} \to \mathcal{K}^n$ is surjective as
$H^n(\mathcal{K}^\bullet) = 0$. As $\mathcal{K}^n$ is flat we see that
$\Ker(\mathcal{K}^{n - 1} \to \mathcal{K}^n)$ is flat
(Modules, Lemma \href{https://stacks.math.columbia.edu/tag/05NK}{lemma flat ses (Tag 05NK)}).
Hence we may replace $\mathcal{K}^\bullet$ by
$\tau_{\leq n - 1}\mathcal{K}^\bullet$. Thus, by induction on $n$, we
reduce to the case that $K^\bullet$ is a complex of flat
$\mathcal{O}_X$-modules with $\mathcal{K}^i = 0$ for $i > b$.

\medskip\noindent
Set $\mathcal{E}^\bullet = \tau_{\geq a}\mathcal{K}^\bullet$.
Everything is clear except that $\mathcal{E}^a$ is flat
which follows immediately from Lemma \ref{cohomology-lemma-last-one-flat}
and the definitions.
\end{proof}


\begin{lemma}
\label{cohomology-lemma-last-one-flat}
Let $(X, \mathcal{O}_X)$ be a ringed space.
Let $\mathcal{E}^\bullet$ be a bounded above complex of flat
$\mathcal{O}_X$-modules with tor-amplitude in $[a, b]$.
Then $\Coker(d_{\mathcal{E}^\bullet}^{a - 1})$ is a flat
$\mathcal{O}_X$-module.
\end{lemma}

\begin{proof}
As $\mathcal{E}^\bullet$ is a bounded above complex of flat modules we see that
$\mathcal{E}^\bullet \otimes_{\mathcal{O}_X} \mathcal{F} =
\mathcal{E}^\bullet \otimes_{\mathcal{O}_X}^{\mathbf{L}} \mathcal{F}$
for any $\mathcal{O}_X$-module $\mathcal{F}$.
Hence for every $\mathcal{O}_X$-module $\mathcal{F}$ the sequence
$$
\mathcal{E}^{a - 2} \otimes_{\mathcal{O}_X} \mathcal{F} \to
\mathcal{E}^{a - 1} \otimes_{\mathcal{O}_X} \mathcal{F} \to
\mathcal{E}^a \otimes_{\mathcal{O}_X} \mathcal{F}
$$
is exact in the middle. Since
$\mathcal{E}^{a - 2} \to \mathcal{E}^{a - 1} \to \mathcal{E}^a \to
\Coker(d^{a - 1}) \to 0$
is a flat resolution this implies that
$\text{Tor}_1^{\mathcal{O}_X}(\Coker(d^{a - 1}), \mathcal{F}) = 0$
for all $\mathcal{O}_X$-modules $\mathcal{F}$. This means that
$\Coker(d^{a - 1})$ is flat, see Lemma \href{https://stacks.math.columbia.edu/tag/08BQ}{lemma flat tor zero (Tag 08BQ)}.
\end{proof}


\begin{lemma}
\label{properties-lemma-finite-locally-free}
Let $X$ be a scheme.
Let $\mathcal{F}$ be a quasi-coherent $\mathcal{O}_X$-module.
The following are equivalent:
\begin{enumerate}
\item $\mathcal{F}$ is a flat $\mathcal{O}_X$-module of finite presentation,
\item $\mathcal{F}$ is $\mathcal{O}_X$-module of finite presentation and
for all $x \in X$ the stalk $\mathcal{F}_x$ is a free
$\mathcal{O}_{X, x}$-module,
\item $\mathcal{F}$ is a locally free, finite type $\mathcal{O}_X$-module,
\item $\mathcal{F}$ is a finite locally free $\mathcal{O}_X$-module, and
\item $\mathcal{F}$ is an $\mathcal{O}_X$-module of finite type,
for every $x \in X$ the stalk $\mathcal{F}_x$ is a free
$\mathcal{O}_{X, x}$-module, and the function
$$
\rho_\mathcal{F} : X \to \mathbf{Z}, \quad
x \longmapsto
\dim_{\kappa(x)} \mathcal{F}_x \otimes_{\mathcal{O}_{X, x}} \kappa(x)
$$
is locally constant in the Zariski topology on $X$.
\end{enumerate}
\end{lemma}

\begin{proof}
This lemma immediately reduces to the affine case.
In this case the lemma is a reformulation of
Algebra, Lemma \ref{algebra-lemma-finite-projective}.
The translation uses
Lemmas \href{https://stacks.math.columbia.edu/tag/01PB}{lemma finite type module (Tag 01PB)},
\href{https://stacks.math.columbia.edu/tag/01PC}{lemma finite presentation module (Tag 01PC)},
\href{https://stacks.math.columbia.edu/tag/05P0}{lemma flat module (Tag 05P0)}, and
\href{https://stacks.math.columbia.edu/tag/05JM}{lemma locally free module (Tag 05JM)}.
\end{proof}


\begin{lemma}
\label{algebra-lemma-finite-projective}
Let $R$ be a ring and let $M$ be an $R$-module.
The following are equivalent
\begin{enumerate}
\item $M$ is finitely presented and $R$-flat,
\item $M$ is finite projective,
\item $M$ is a direct summand of a finite free $R$-module,
\item $M$ is finitely presented and
for all $\mathfrak p \in \Spec(R)$ the
localization $M_{\mathfrak p}$ is free,
\item $M$ is finitely presented and
for all maximal ideals $\mathfrak m \subset R$ the
localization $M_{\mathfrak m}$ is free,
\item $M$ is finite and locally free,
\item $M$ is finite locally free, and
\item $M$ is finite, for every prime $\mathfrak p$ the module
$M_{\mathfrak p}$ is free, and the function
$$
\rho_M : \Spec(R) \to \mathbf{Z}, \quad
\mathfrak p
\longmapsto
\dim_{\kappa(\mathfrak p)} M \otimes_R \kappa(\mathfrak p)
$$
is locally constant in the Zariski topology.
\end{enumerate}
\end{lemma}

\begin{proof}
First suppose $M$ is finite projective, i.e., (2) holds.
Take a surjection $R^n \to M$ and let $K$ be the kernel.
Since $M$ is projective,
$0 \to K \to R^n \to M \to 0$ splits.
Hence (2) $\Rightarrow$ (3).
The implication (3) $\Rightarrow$ (2) follows from the fact that
a direct summand of a projective is projective, see
Lemma \href{https://stacks.math.columbia.edu/tag/05CF}{lemma characterize projective (Tag 05CF)}.

\medskip\noindent
Assume (3), so we can write $K \oplus M \cong R^{\oplus n}$.
So $K$ is a  direct summand of $R^n$ and thus finitely generated.
This shows $M = R^{\oplus n}/K$ is finitely presented.
In other words, (3) $\Rightarrow$ (1).

\medskip\noindent
Assume $M$ is finitely presented and flat, i.e., (1) holds.
We will prove that (7) holds. Pick any prime $\mathfrak p$ and
$x_1, \ldots, x_r \in M$ which map to a basis of
$M \otimes_R \kappa(\mathfrak p)$. By
Nakayama's lemma (in the form of Lemma \href{https://stacks.math.columbia.edu/tag/0GLX}{lemma NAK localization (Tag 0GLX)})
these elements generate $M_g$ for some $g \in R$, $g \not \in \mathfrak p$.
The corresponding surjection $\varphi : R_g^{\oplus r} \to M_g$
has the following two properties: (a) $\Ker(\varphi)$ is a finite
$R_g$-module (see Lemma \href{https://stacks.math.columbia.edu/tag/0519}{lemma extension (Tag 0519)})
and (b) $\Ker(\varphi) \otimes \kappa(\mathfrak p) = 0$
by flatness of $M_g$ over $R_g$ (see
Lemma \href{https://stacks.math.columbia.edu/tag/00HL}{lemma flat tor zero (Tag 00HL)}).
Hence by Nakayama's lemma again there exists a $g' \in R_g$ such that
$\Ker(\varphi)_{g'} = 0$. In other words, $M_{gg'}$ is free.

\medskip\noindent
A finite locally free module is a finite module, see
Lemma \href{https://stacks.math.columbia.edu/tag/00EO}{lemma cover (Tag 00EO)},
hence (7) $\Rightarrow$ (6).
It is clear that (6) $\Rightarrow$ (7) and that (7) $\Rightarrow$ (8).

\medskip\noindent
A finite locally free module is a finitely presented module, see
Lemma \href{https://stacks.math.columbia.edu/tag/00EO}{lemma cover (Tag 00EO)},
hence (7) $\Rightarrow$ (4).
Of course (4) implies (5).
Since we may check flatness locally (see
Lemma \href{https://stacks.math.columbia.edu/tag/00HT}{lemma flat localization (Tag 00HT)})
we conclude that (5) implies (1).
At this point we have
$$
\xymatrix{
(2) \ar@{<=>}[r] & (3) \ar@{=>}[r] & (1) \ar@{=>}[r] &
(7)  \ar@{<=>}[r] \ar@{=>}[rd] \ar@{=>}[d] & (6) \\
& & (5) \ar@{=>}[u] & (4) \ar@{=>}[l] & (8)
}
$$

\medskip\noindent
Suppose that $M$ satisfies (1), (4), (5), (6), and (7).
We will prove that (3) holds. It suffices
to show that $M$ is projective. We have to show that $\Hom_R(M, -)$
is exact. Let $0 \to N'' \to N \to N'\to 0$ be a short exact sequence of
$R$-module. We have to show that
$0 \to \Hom_R(M, N'') \to \Hom_R(M, N) \to
\Hom_R(M, N') \to 0$ is exact.
As $M$ is finite locally free there exist a covering
$\Spec(R) = \bigcup D(f_i)$ such that $M_{f_i}$ is finite free.
By
Lemma \href{https://stacks.math.columbia.edu/tag/0583}{lemma hom from finitely presented (Tag 0583)}
we see that
$$
0 \to \Hom_R(M, N'')_{f_i} \to \Hom_R(M, N)_{f_i} \to
\Hom_R(M, N')_{f_i} \to 0
$$
is equal to
$0 \to \Hom_{R_{f_i}}(M_{f_i}, N''_{f_i}) \to
\Hom_{R_{f_i}}(M_{f_i}, N_{f_i}) \to
\Hom_{R_{f_i}}(M_{f_i}, N'_{f_i}) \to 0$
which is exact as $M_{f_i}$ is free and as the localization
$0 \to N''_{f_i} \to N_{f_i} \to N'_{f_i} \to 0$
is exact (as localization is exact). Whence we see that
$0 \to \Hom_R(M, N'') \to \Hom_R(M, N) \to
\Hom_R(M, N') \to 0$ is exact by
Lemma \href{https://stacks.math.columbia.edu/tag/00EO}{lemma cover (Tag 00EO)}.

\medskip\noindent
Finally, assume that (8) holds. Pick a maximal ideal $\mathfrak m \subset R$.
Pick $x_1, \ldots, x_r \in M$ which map to a $\kappa(\mathfrak m)$-basis of
$M \otimes_R \kappa(\mathfrak m) = M/\mathfrak mM$. In particular
$\rho_M(\mathfrak m) = r$. By
Nakayama's Lemma \href{https://stacks.math.columbia.edu/tag/00DV}{lemma NAK (Tag 00DV)}
there exists an $f \in R$, $f \not \in \mathfrak m$ such that
$x_1, \ldots, x_r$ generate $M_f$ over $R_f$. By the assumption that
$\rho_M$ is locally constant there exists a $g \in R$, $g \not \in \mathfrak m$
such that $\rho_M$ is constant equal to $r$ on $D(g)$. We claim that
$$
\Psi : R_{fg}^{\oplus r} \longrightarrow M_{fg}, \quad
(a_1, \ldots, a_r) \longmapsto \sum a_i x_i
$$
is an isomorphism. This claim will show that $M$ is finite locally
free, i.e., that (7) holds. To see the claim
it suffices to show that the induced map on localizations
$\Psi_{\mathfrak p} : R_{\mathfrak p}^{\oplus r} \to M_{\mathfrak p}$
is an isomorphism for all $\mathfrak p \in D(fg)$, see
Lemma \href{https://stacks.math.columbia.edu/tag/00HN}{lemma characterize zero local (Tag 00HN)}.
By our choice of $f$ the map $\Psi_{\mathfrak p}$
is surjective. By assumption (8) we have
$M_{\mathfrak p} \cong R_{\mathfrak p}^{\oplus \rho_M(\mathfrak p)}$
and by our choice of $g$ we have $\rho_M(\mathfrak p) = r$.
Hence $\Psi_{\mathfrak p}$ determines a surjection
$R_{\mathfrak p}^{\oplus r} \to
M_{\mathfrak p} \cong R_{\mathfrak p}^{\oplus r}$
whence is an isomorphism by
Lemma \href{https://stacks.math.columbia.edu/tag/05G8}{lemma fun (Tag 05G8)}.
(Of course this last fact follows from a simple matrix argument also.)
\end{proof}


\begin{lemma}
\label{derived-lemma-cohomology-homological}
Let $\mathcal{A}$ be an abelian category. The functor
$$
H^0 : K(\mathcal{A}) \longrightarrow \mathcal{A}
$$
is homological.
\end{lemma}

\begin{proof}
Because $H^0$ is a functor, and by our definition of distinguished triangles
it suffices to prove that given a termwise split short exact sequence
of complexes $0 \to A^\bullet \to B^\bullet \to C^\bullet \to 0$
the sequence $H^0(A^\bullet) \to H^0(B^\bullet) \to H^0(C^\bullet)$
is exact. This follows from
Homology, Lemma \href{https://stacks.math.columbia.edu/tag/0117}{lemma long exact sequence cochain (Tag 0117)}.
\end{proof}


\begin{lemma}
\label{cohomology-lemma-cone-tor-amplitude}
Let $(X, \mathcal{O}_X)$ be a ringed space.
Let $(K, L, M, f, g, h)$ be a distinguished
triangle in $D(\mathcal{O}_X)$. Let $a, b \in \mathbf{Z}$.
\begin{enumerate}
\item If $K$ has tor-amplitude in $[a + 1, b + 1]$ and
$L$ has tor-amplitude in $[a, b]$ then $M$ has
tor-amplitude in $[a, b]$.
\item If $K$ and $M$ have tor-amplitude in $[a, b]$, then
$L$ has tor-amplitude in $[a, b]$.
\item If $L$ has tor-amplitude in $[a + 1, b + 1]$
and $M$ has tor-amplitude in $[a, b]$, then
$K$ has tor-amplitude in $[a + 1, b + 1]$.
\end{enumerate}
\end{lemma}

\begin{proof}
Omitted. Hint: This just follows from the long exact cohomology sequence
associated to a distinguished triangle and the fact that
$- \otimes_{\mathcal{O}_X}^{\mathbf{L}} \mathcal{F}$
preserves distinguished triangles.
The easiest one to prove is (2) and the others follow from it by
translation.
\end{proof}



\noindent
Let $(X, \mathcal{O}_X)$ be a ringed space.
Let $\mathcal{F}^\bullet$ be an object of $D(\mathcal{O}_X)$.
Choose a K-flat resolution $\mathcal{K}^\bullet \to \mathcal{F}^\bullet$, see
Lemma \href{https://stacks.math.columbia.edu/tag/06YF}{lemma K flat resolution (Tag 06YF)}.
By
Lemma \href{https://stacks.math.columbia.edu/tag/06Y8}{lemma derived tor exact (Tag 06Y8)}
we obtain an exact functor of triangulated categories
$$
K(\mathcal{O}_X)
\longrightarrow
K(\mathcal{O}_X),
\quad
\mathcal{G}^\bullet
\longmapsto
\text{Tot}(\mathcal{G}^\bullet \otimes_{\mathcal{O}_X} \mathcal{K}^\bullet)
$$
By
Lemma \href{https://stacks.math.columbia.edu/tag/06YA}{lemma K flat quasi isomorphism (Tag 06YA)}
this functor induces a functor
$D(\mathcal{O}_X) \to D(\mathcal{O}_X)$ simply because
$D(\mathcal{O}_X)$ is the localization of $K(\mathcal{O}_X)$
at quasi-isomorphisms. As the category of $K$-flat resolutions
of $\mathcal{F}^\bullet$ is cofiltered and as we have
Lemma \href{https://stacks.math.columbia.edu/tag/06YG}{lemma derived tor quasi isomorphism other side (Tag 06YG)}
the resulting functor (up to isomorphism)
does not depend on the choice of $\mathcal{K}^\bullet$.

\begin{definition}
\label{cohomology-definition-derived-tor}
Let $(X, \mathcal{O}_X)$ be a ringed space.
Let $\mathcal{F}^\bullet$ be an object of $D(\mathcal{O}_X)$.
The {\it derived tensor product}
$$
- \otimes_{\mathcal{O}_X}^{\mathbf{L}} \mathcal{F}^\bullet :
D(\mathcal{O}_X)
\longrightarrow
D(\mathcal{O}_X)
$$
is the exact functor of triangulated categories described above.
\end{definition}

\noindent
It is clear from our explicit constructions that
there is a canonical isomorphism
$$
\mathcal{F}^\bullet \otimes_{\mathcal{O}_X}^{\mathbf{L}} \mathcal{G}^\bullet
\cong
\mathcal{G}^\bullet \otimes_{\mathcal{O}_X}^{\mathbf{L}} \mathcal{F}^\bullet
$$
for $\mathcal{G}^\bullet$ and $\mathcal{F}^\bullet$ in $D(\mathcal{O}_X)$.


\noindent
In particular, this lemma implies that a distinguished triangle
$(X, Y, Z, f, g, h)$ in $K(\mathcal{A})$ gives rise to a long exact
cohomology sequence
\begin{equation}
\label{derived-equation-long-exact-cohomology-sequence-D}
\xymatrix{
\ldots \ar[r] &
H^i(X) \ar[r]^{H^i(f)} &
H^i(Y) \ar[r]^{H^i(g)} &
H^i(Z) \ar[r]^{H^i(h)} &
H^{i + 1}(X) \ar[r] & \ldots
}
\end{equation}
see (\href{https://stacks.math.columbia.edu/tag/0148}{equation long exact cohomology sequence (Tag 0148)}). Moreover, there is
a compatibility with the long exact sequence of cohomology associated to
a short exact sequence of complexes (insert future reference here). For
example, if $(A^\bullet, B^\bullet, C^\bullet, \alpha, \beta, \delta)$
is the distinguished triangle associated to a termwise split exact
sequence of complexes (see
Definition \ref{derived-definition-split-ses}),
then the cohomology sequence above agrees with the one defined using the
snake lemma, see
Homology, Lemma \href{https://stacks.math.columbia.edu/tag/0117}{lemma long exact sequence cochain (Tag 0117)}
and for agreement of sequences, see
Homology, Lemma \href{https://stacks.math.columbia.edu/tag/011K}{lemma ses termwise split long cochain (Tag 011K)}.

\medskip\noindent
Recall that a complex $K^\bullet$ is {\it acyclic} if $H^i(K^\bullet) = 0$
for all $i \in \mathbf{Z}$. Moreover, recall that a morphism of complexes
$f : K^\bullet \to L^\bullet$ is a {\it quasi-isomorphism} if and only if
$H^i(f)$ is an isomorphism for all $i$. See
Homology, Definition \ref{homology-definition-quasi-isomorphism-cochain}.



\begin{lemma}
\label{lemma-construct-strictly-perfect}
Let $X$ be a quasi-compact and quasi-separated scheme with the
resolution property.
Let $\mathcal{F}^\bullet$ be a bounded below complex of quasi-coherent
$\mathcal{O}_X$-modules representing a perfect object of
$D(\mathcal{O}_X)$. Then there exists a bounded complex
$\mathcal{E}^\bullet$ of finite locally free $\mathcal{O}_X$-modules
and a quasi-isomorphism $\mathcal{E}^\bullet \to \mathcal{F}^\bullet$.
\end{lemma}

\begin{proof}
Let $a, b \in \mathbf{Z}$ be integers such that $\mathcal{F}^\bullet$
has tor amplitude in $[a, b]$ and such that $\mathcal{F}^n = 0$ for
$n < a$. The existence of such a pair of integers
follows from Cohomology, Lemma \ref{cohomology-lemma-perfect}
and the fact that $X$ is quasi-compact.
If $b < a$, then $\mathcal{F}^\bullet$ is zero in the
derived category and the lemma holds.
We will prove by induction on
$b - a \geq 0$ that there exists a complex
$\mathcal{E}^a \to \ldots \to \mathcal{E}^b$
with $\mathcal{E}^i$ finite locally free
and a quasi-isomorphism $\mathcal{E}^\bullet \to \mathcal{F}^\bullet$.

\medskip\noindent
The base case is the case $b - a = 0$. In this case
$H^b(\mathcal{F}^\bullet) = H^a(\mathcal{F}^\bullet) =
\Ker(\mathcal{F}^a \to \mathcal{F}^{a + 1})$
is finite locally free. Namely, it is a finitely presented
$\mathcal{O}_X$-module of tor dimension $0$ and hence finite
locally free. See Cohomology, Lemmas \ref{cohomology-lemma-perfect} and
\ref{cohomology-lemma-finite-cohomology} and
Properties, Lemma \ref{properties-lemma-finite-locally-free}.
Thus we can take $\mathcal{E}^\bullet$ to be
$H^b(\mathcal{F}^\bullet)$ sitting in degree $b$.
The rest of the proof is dedicated to the induction step.

\medskip\noindent
Assume $b > a$. Observe that
$$
H^b(\mathcal{F}^\bullet) =
\Ker(\mathcal{F}^b \to \mathcal{F}^{b + 1})/
\Im(\mathcal{F}^{b - 1} \to \mathcal{F}^b)
$$
is a finite type quasi-coherent $\mathcal{O}_X$-module, see
Cohomology, Lemmas \ref{cohomology-lemma-perfect} and
\ref{cohomology-lemma-finite-cohomology}. Then we can find a finite type
quasi-coherent $\mathcal{O}_X$-module $\mathcal{F}$ and a map
$$
\mathcal{F} \longrightarrow
\Ker(\mathcal{F}^b \to \mathcal{F}^{b + 1})
$$
such that the composition with the projection onto
$H^b(\mathcal{F}^\bullet)$ is surjective.
Namely, we can write $\Ker(\mathcal{F}^b \to \mathcal{F}^{b + 1})$
as the filtered union of its finite type quasi-coherent submodules by
Properties, Lemma \ref{properties-lemma-quasi-coherent-colimit-finite-type}
and then one of these will do the job.
Next, we choose a finite locally free $\mathcal{O}_X$-module
$\mathcal{E}^b$ and a surjection $\mathcal{E}^b \to \mathcal{F}$ using
the resolution property of $X$. Consider the map of complexes
$$
\alpha : \mathcal{E}^b[-b] \to \mathcal{F}^\bullet
$$
and its cone $C(\alpha)^\bullet$, see
Derived Categories, Definition \ref{derived-definition-cone}.
We observe that $C(\alpha)^\bullet$ is nonzero only in degrees
$\geq a$, has tor amplitude in $[a, b]$ by
Cohomology, Lemma \ref{cohomology-lemma-cone-tor-amplitude},
and has $H^b(C(\alpha)^\bullet) = 0$ by construction.
Thus we actually find that $C(\alpha)^\bullet$ has tor amplitude
in $[a, b - 1]$. Hence the induction hypothesis applies to
$C(\alpha)^\bullet$ and we find a map of complexes
$$
(\mathcal{E}^a \to \ldots \to \mathcal{E}^{b - 1})
\longrightarrow
C(\alpha)^\bullet
$$
with properties as stated in the induction hypothesis. Unwinding
the definition of the cone this gives a commutative diagram
$$
\xymatrix{
\ldots \ar[r] &
\mathcal{E}^{b - 2} \ar[r] \ar[d] &
\mathcal{E}^{b - 1} \ar[r] \ar[d] &
0 \ar[r] \ar[d] &
\ldots \\
\ldots \ar[r] &
\mathcal{F}^{b - 2} \ar[r] &
\mathcal{F}^{b - 1} \oplus \mathcal{E}^b \ar[r] &
\mathcal{F}^b \ar[r] &
\ldots
}
$$
It is clear that we obtain a map of complexes
$(\mathcal{E}^a \to \ldots \to \mathcal{E}^b) \to \mathcal{F}^\bullet$.
We omit the verification that this map is a quasi-isomorphism.
\end{proof}
\end{document}
